\pdfoutput=1
\documentclass[11pt]{article}
\usepackage[margin=25mm]{geometry}
\usepackage{amsmath,amssymb,amsthm,mathtools}
\usepackage[hidelinks,pdftitle={Regge Refinement and the Vacuum High-Frequency Limit Do Not Commute}]{hyperref}
\newtheorem{theorem}{Theorem}
\newtheorem{lemma}[theorem]{Lemma}
\newtheorem{corollary}[theorem]{Corollary}
\DeclareMathOperator{\Ric}{Ric}
\DeclareMathOperator{\tr}{tr}
\newcommand{\dd}{\mathrm d}
\newcommand{\Rcal}{\mathcal R}
\newcommand{\Ecal}{\mathcal E}
\newcommand{\Hdual}{(H^2_0)^*}

\title{Regge Refinement and the\\Vacuum High-Frequency Limit Do Not Commute}
\author{}
\date{}

\begin{document}
\maketitle

\begin{abstract}
A single family of admissible tilted Lorentzian simplex meshes gives
different limits of the full signed weak Einstein density of
interpolated exact vacuum plane waves. Refining the affine
tangential-moment Regge interpolants first yields zero. Sending the
wavelength to zero first yields the nonzero null-dust Einstein
density of the smooth limiting Rosen metric. We track the background
derivative dependence of the signed bulk, face and joint variations
to prove the same alternatives on simultaneous resolved and
underresolved diagonals. At one comparable-scale resonance of a
rationally tilted mesh, all affine moments of the leading cosine
quadrature vanish; the sine quadrature does not disappear. Two
explicit exact circularly polarized waves with opposite continuum
helicities then become indistinguishable to bounded quadratic
broken-gradient readouts on that diagonal, and their signed weak Einstein
functionals differ by only \(O(h)\). The results concern the
variational residual of degree-one interpolated metrics, not the
convergence of solutions of a discrete field equation.
\end{abstract}

\section{Two independent limits}

In general, a uniform limit of exact vacuum metrics need not satisfy
the vacuum equation. Burnett's weak-limit formulation identifies a
quadratic derivative defect, and the work of Green and Wald gives
explicit continuum examples and restrictions on the resulting
effective stress \cite{Burnett,GreenWald}. More recently Huneau and
Luk prove a massless-Vlasov description for the stated class of
generalized-wave-coordinate limits \cite{HuneauLuk}. None of those
continuum statements compares a wave family with a particular
signed Lorentzian Regge interpolation at finite mesh size.
Conversely a fixed-background Regge consistency estimate may have
a constant that grows with the derivative norms of the metric.

We make the comparison on a single common family of meshes. In a
bounded coordinate domain \(\Omega\) in four dimensions, define
the smooth Einstein-density functional, for covariant
symmetric test tensors, by
\begin{equation}
 \Ecal(g)(\sigma)=\int_\Omega G(g)^{ab}\sigma_{ab}
                            \sqrt{|\det g|}\,\dd^4x .
 \label{eq:smooth}
\end{equation}
The indices on \(G=\Ric-\tfrac12 Rg\) are raised using \(g\).
The norm \(H^{-2}\) below means the dual of
\(H^2_0(\Omega;\operatorname{Sym}^2)\) for a fixed auxiliary Euclidean
coordinate norm, not a Lorentzian spacetime Sobolev norm.

We use the signed real Lorentzian action derivative of
\cite{SignedRegge}, abbreviated \(\Rcal_h(m)\), with the same
curvature, reference-subtraction and face orientation conventions.
More explicitly, on an admissible piecewise smooth Regge metric \(m\),
\begin{align}
 \Rcal_h(m)(\sigma)
 ={}&\sum_K\int_K G(m)^{ab}\sigma_{ab}\,\dd\mu_m \notag\\
 &-\sum_F \epsilon_F\int_F
      \langle[K_F]-[H_F]m_F,\sigma_F\rangle_{m_F}\,\dd\mu_F
   -\sum_J\int_J k_J\tr_{m_J}(\sigma_J)\,\dd\mu_J .
 \label{eq:regge}
\end{align}
Here \(F,J\) denote interior faces and hinges, \(\epsilon_F\)
is the squared length of the unit cooriented normal, \(K_F,H_F\)
are the second fundamental form and its trace, and brackets are jumps
in that coorientation. The real angle deficit \(k_J\) is on the
regular branch continued from the smooth zero-deficit fan; its
turning reference has already been subtracted. Face and hinge
test contractions use the induced forms, and all the measures
are the actual absolute induced measures. For smooth \(m\) the
interior correction terms vanish. The convention for real
non-null joints and angle branches is consistent with
\cite{Hayward,Sorkin,Asante}.

On each simplex let \(I_h\) be the degree-one tangential-moment
interpolant of \cite{SignedRegge}, specified by degree-zero and
degree-one edge moments and the constant tangential tensor moments
on each triangular face. In particular \(I_h\) is not the
edge-length-only interpolation in four dimensions. A useful
Riemannian precursor for variational Einstein tensors is
\cite{GawlikNeunteufel}; the present comparison uses the signed
Lorentzian form of \eqref{eq:regge}.

\begin{theorem}[Two noncommuting limits]\label{thm:limits}
There exist a smooth exact vacuum family \(g_\varepsilon\) converging
uniformly to a smooth metric \(g_0\), and a single regular, non-null,
shape-regular quasi-uniform family of tilted affine four-simplex
meshes admissible for all the comparisons, such that in \(H^{-2}\)
\begin{equation}
 \lim_{\varepsilon\downarrow0}\lim_{h\downarrow0}
       \Rcal_h(I_hg_\varepsilon)=0,\qquad
 \lim_{h\downarrow0}\lim_{\varepsilon\downarrow0}
       \Rcal_h(I_hg_\varepsilon)=\Ecal(g_0)\ne0 .
 \label{eq:noncommute}
\end{equation}
There is a common constant \(C\) and a fixed
\(\varepsilon_0>0\) for which
\begin{align}
 \|\Rcal_h(I_hg_\varepsilon)\|_{H^{-2}}
   &\le C(h/\varepsilon)^2 ,
               && 0<h\le\varepsilon\le\varepsilon_0,\label{eq:res}\\
 \|\Rcal_h(I_hg_\varepsilon)-\Ecal(g_0)\|_{H^{-2}}
   &\le C\bigl[h^2+(\varepsilon/h)^2\bigr].
               &&\label{eq:under}
\end{align}
In particular the corresponding simultaneous limits are zero
for \(h=o(\varepsilon)\) and nonzero for \(\varepsilon=o(h)\).
The nonzero density can have any sufficiently small smooth
square envelope along the chosen Rosen null coordinate.
\end{theorem}

Sections \ref{sec:waves}--\ref{sec:bounds} prove this theorem.
Sections \ref{sec:alias} and \ref{sec:helicity} then analyze a single
fixed comparable-scale resonance, without asserting a universal
intermediate scaling law.

\section{Exact vacuum waves and their limit}\label{sec:waves}

On a bounded Rosen chart let
\begin{equation}
 g_\varepsilon=-2\,\dd u\,\dd v
     +a_\varepsilon(u)^2
       \left(e^{2w_\varepsilon(u)}\dd x^2
                      +e^{-2w_\varepsilon(u)}\dd y^2\right),\qquad
 w_\varepsilon(u)=\varepsilon A(u)
                           \sin(\phi(u)/\varepsilon).
 \label{eq:wave}
\end{equation}
The smooth phase \(\phi\) has \(\phi'>0\) bounded away from zero
on a fixed larger chart. Scale a fixed smooth nonzero envelope
\(A\) small enough that its amplitude and derivatives are small
on this chart. Put \(B=A\phi'\). At a fixed \(u_*\) impose
\(a_\varepsilon=1,\ a_\varepsilon'=0\) and solve
\begin{equation}
        a_\varepsilon''+(w_\varepsilon')^2a_\varepsilon=0.
 \label{eq:vac}
\end{equation}
The ODE gives uniform positivity of \(a_\varepsilon\) on the
bounded chart for sufficiently small amplitude. For transverse
factors \(p=a_\varepsilon e^{w_\varepsilon}\) and
\(q=a_\varepsilon e^{-w_\varepsilon}\), the Rosen Ricci formula
has only possible nonzero component
\[
       \Ric_{uu}=-p''/p-q''/q
                   =-2\big(a_\varepsilon''/a_\varepsilon
                                    +(w_\varepsilon')^2\big).
\]
Thus \eqref{eq:wave} is exactly vacuum, not a truncated perturbation.

The identity
\[
 (w_\varepsilon')^2
    =\tfrac12 B^2(1+\cos(2\phi/\varepsilon))+O(\varepsilon)
\]
gives a uniform bound on its coefficient and an \(O(\varepsilon)\)
primitive for its difference from \(B^2/2\): integrate the cosine
once by parts using \(\phi'\ne0\). Let \(a_0\) solve
\(a_0''+\frac12B^2a_0=0\) with the same initial data.
The difference equation in Volterra form has an
\(O(\varepsilon)\) integrated forcing, also when its first
derivative is taken. Indeed apply the same integration by
parts with the uniformly \(C^1\)-bounded factor \(a_\varepsilon\)
from \eqref{eq:vac}. Gronwall yields
\(\|a_\varepsilon-a_0\|_{C^1}\lesssim\varepsilon\). Consequently
\begin{align}
 g_0&=-2\,\dd u\,\dd v+a_0^2(\dd x^2+\dd y^2),&
 \|g_\varepsilon-g_0\|_\infty&\lesssim\varepsilon,\notag\\
 \|Dg_\varepsilon\|_\infty&\lesssim1,&
 \|D^2 g_\varepsilon\|_\infty&\lesssim\varepsilon^{-1}.
 \label{eq:jets}
\end{align}
The metric \(g_0\) has \(\Ric_{uu}=B^2\) and scalar curvature zero.
With the raised-index convention in \eqref{eq:smooth}, its
nonzero density is
\begin{equation}
  \Ecal(g_0)(\sigma)=\int_\Omega
            B(u)^2a_0(u)^2\,\sigma_{vv}\,\dd u\,\dd v\,\dd x\,\dd y.
 \label{eq:dust}
\end{equation}
Thus a compactly supported nonnegative \(\sigma_{vv}\) tests the
defect strictly on any open set where \(B\ne0\). Choosing \(\phi=u\)
and prescribing small \(A=B\) gives the stated square envelope.
For \(A=\kappa_0>0\), \(\phi=u\) and \(u_*=0\), one has
\(a_0=\cos(\kappa_0 u/\sqrt2)\).

The normalization can also be checked from the rapidly varying
shear. In the fixed limiting transverse frame
\[
 \partial_u(g_\varepsilon-g_0)_{AB}
       =2a_0^2 B\cos(\phi/\varepsilon)
                   \operatorname{diag}(1,-1)_{AB}+O(\varepsilon).
\]
The derivative converges weakly to zero, whereas one quarter of
its transverse squared norm (using \(g_0\) to contract the screen)
converges weakly to \(B^2\). This computes the radiative defect
directly on the explicit family. It uses no generalized-wave
coordinate condition for the interpolants. The effective physical
stress is \(B^2\dd u\otimes\dd u/(8\pi G_N)\) if the Einstein
normalization is restored.

\section{A common non-null tilted mesh}\label{sec:mesh}

Use a reference Minkowski chart related rationally to the null
chart by \(u=t+z,\ v=(t-z)/2\). A finite affine subdivision of a
periodic box template has finitely many reference face and hinge
tangent subspaces. Choose a fixed arbitrarily small generic linear
tilt \(L\). Vanishing of any normalized face or hinge tangent
Gram determinant is a proper polynomial condition in \(L\):
one can send that fixed subspace to a nondegenerate plane by
a linear isomorphism. We also exclude the finitely many linear
conditions \(\dd u(L e)=0\) on template edge directions, needed
only for the resonance argument below. A rational tilt exists
outside this finite union. Some template faces and hinges can
be chosen strictly spacelike and others strictly timelike before
tilting; a small tilt preserves these strict signs. For example,
one may use a rectangular triangulation with the temporal axis
shorter than the spatial axes. A fixed tilted parallelepiped or
union of boxes is exactly filled by uniform subdivisions of the
same template. We call its lattice parameter \(h\), uniformly
equivalent to the maximum simplex diameter, and restrict tests
to the fixed bounded interior chart.

All normalized reference Gram margins and all incidence and
shape constants are uniform in \(h\). For the moment scheme
\(I_h\), reference-simplex unisolvence and affine scaling give
local \(L^\infty\) stability, reproduction of affine tensors,
and the affine inverse estimate
\begin{equation}
 \|I_hs\|_{\infty,K}\le C_I\|s\|_{\infty,K},\qquad
 \|D I_hs\|_{\infty,K}\le C_Ih^{-1}\|s\|_{\infty,K}.
 \label{eq:interp}
\end{equation}
For precision, edge restrictions of an affine tensor are fixed
by its two moments. On an affine four-simplex the ten
\(d\lambda_i\odot d\lambda_j\), \(i<j\), form a tensor basis
adapted to these restrictions. Once edge moments vanish, the
coefficient of each basis tensor is a linear combination of
barycentric coordinates of the other vertices. Each triangular
tangential test fixes its corresponding remaining coefficients
by integration of a single complementary barycentric coordinate.
This gives local unisolvence and identical tangential face
restrictions on adjoining cells with shared moments
\cite{SignedRegge}. No positive-definiteness of the input tensor
is needed for this linear calculation.

The exact metrics above, their limits and the later examples
can all be made sufficiently close to the same constant
\(\eta\). By \eqref{eq:interp} and \(I_h\eta=\eta\),
their interpolants remain uniformly close to \(\eta\) too.
Choose the amplitudes and the chart so that a common admissible
\(C^0\) ball contains both original metrics and all piecewise
straight comparison paths. Every point on such a path has
a common strict timelike covector, nondegenerate tangential
forms and bounded inverse metric. This step requires no
uniform higher-derivative assumption.

At a reference spacelike hinge the Lorentzian normal plane
has a genuine cyclic fan with precisely four null rays and
no null bounding rays. A bounding ray's signed squared
length is a Schur complement of the hinge block in the
induced face form. It keeps its causal type within the
chosen ball; that type is the same on both sides of a
shared boundary by tangential continuity. The orientation
and the light-ray count in each sector cannot change without
a degenerate sector or a null bounding ray. Thus all paths
lie on the real reference branch of \eqref{eq:regge}, with
uniform margins for its angle coefficients and their metric
derivatives up to any fixed order needed below. These statements
are local to this admissible template, not a claim about
arbitrary causality assignments on Lorentzian Regge meshes.

\section{A two-scale bound for the full weak residual}\label{sec:bounds}

Let \(g\) be smooth and \(m=g+e\) be tangentially continuous and
affine on the common mesh, with the straight comparison path
in the fixed admissible ball. Let \(A_j=\|D^j e\|_{\infty,h}\),
\(j=0,1,2\), in broken norms with one-sided traces, and
\(B_2=\|D^2g\|_\infty\); \(\|Dg\|_\infty\) is uniformly bounded
for the backgrounds below. The estimate useful here is
\begin{align}
 \|\Rcal_h(g+e)-\Ecal(g)\|_{H^{-2}}\le C\{&
 A_0(1+B_2)+(1+B_2)A_0^2+A_0A_1+A_1^2+A_0A_2
 \notag\\[-2pt]
 &+h^{-1}(A_0^2+A_0A_1)+h^{-2}A_0^2\},
 \label{eq:uniform}
\end{align}
with \(C\) depending on the common margins, shape family,
domain and first background jet bound only.

To obtain the linear term, commute mixed variations of the
signed action and vary the test on the smooth metric side,
as in \cite{SignedRegge}. At a smooth metric the face and hinge
terms cancel, so for a compact smooth test
\begin{equation}
 D\Rcal_h(g)[e](\sigma)
   =\int_\Omega e_{ab}\,
        D\!\left(\sqrt{|\det g|}\,G(g)^{ab}\right)[\sigma]\,\dd^4x .
 \label{eq:hess}
\end{equation}
External boundary terms may be included, or \(e\) cut off away
from the support of \(\sigma\). This identity transfers
the complete first-order sum, not its separate singular
pieces. The smooth differential operator on \(\sigma\) is
of order at most two. Its coefficients use only bounded
functions of \(g,g^{-1}\), first jets up to degree two
and the background second jet linearly. Thus the linear
bound is \(CA_0(1+B_2)\|\sigma\|_{H^2}\).

For the nonlinear remainder, each bulk density is linear
in the second metric jet and quadratic in the first metric
jet, with uniformly smooth coefficients in the metric values.
Taylor expansion after removing the linear term gives
\[
 C\big((1+B_2)A_0^2+A_0A_1+A_1^2+A_0A_2\big)|\sigma|.
\]
In particular no term \(A_2^2\) can arise. On an affine face
the traction, including the density and test contractions, is
linear in the first one-sided metric jets; its remainder is
\(C(A_0^2+A_0A_1)|\sigma_F|\) per unit auxiliary face area.
The hinge bracket depends only on metric values within its
fixed branch, giving \(CA_0^2|\sigma_J|\) per unit auxiliary
hinge area. In these expressions the actual signs, densities
and inverse contractions are in the smooth coefficients.
The scaled trace bounds on the auxiliary affine simplexes
give
\[
 \sum_F\|\sigma\|_{L^2(F)}^2\lesssim h^{-1}\|\sigma\|_{H^2}^2,
 \qquad
 \sum_J\|\sigma\|_{L^2(J)}^2\lesssim h^{-2}\|\sigma\|_{H^2}^2 .
\]
Indeed the codimension-two trace is controlled by \(H^2\) on
the fixed reference simplex and then scaled, and each cell
contributes boundedly many strata. Total face and hinge areas
are \(O(h^{-1})\) and \(O(h^{-2})\). Cauchy--Schwarz gives the
corresponding sum factors in \eqref{eq:uniform}. Combining
with \eqref{eq:hess} proves \eqref{eq:uniform} and extends it
to \(H^2_0\) by density.

\begin{proof}[Proof of Theorem~\ref{thm:limits}]
For the resolved comparison use \(g=g_\varepsilon\). Subtracting
the linear Taylor polynomial on each cell in \eqref{eq:interp}
and using \eqref{eq:jets}, we have
\(A_0\lesssim h^2/\varepsilon\), \(A_1\lesssim h/\varepsilon\),
\(A_2,B_2\lesssim\varepsilon^{-1}\).
When \(h\le\varepsilon\le1\), every term of
\eqref{eq:uniform} is \(O(h^2/\varepsilon^2)\), including
\(B_2A_0^2=O(h^4/\varepsilon^3)\).
Since \(\Ecal(g_\varepsilon)=0\), this is \eqref{eq:res}.

For the underresolved comparison use \(g=g_0\) and
\[
       e=(I_hg_0-g_0)+I_h(g_\varepsilon-g_0).
\]
By \eqref{eq:interp} and \eqref{eq:jets},
\(A_0\lesssim h^2+\varepsilon\), \(A_1\lesssim h+\varepsilon/h\),
and \(A_2\lesssim1\) because \(I_hg_\varepsilon\) is exactly
affine. Here \(B_2\) is fixed. Equation \eqref{eq:uniform}
gives \(C(h^2+\varepsilon+(\varepsilon/h)^2)\), implying
\eqref{eq:under} by \(2\varepsilon\le h^2+(\varepsilon/h)^2\).
This avoids inserting the large second jet of \(g_\varepsilon\)
in the underresolved comparison.

At fixed \(h\), finite-dimensional interpolation and
\(g_\varepsilon\to g_0\) uniformly imply convergence of all
broken affine coefficients and first derivatives.
On the fixed admissible branch, the finitely many bulk,
face and hinge coefficients of \eqref{eq:regge} then converge
in the \(H^{-2}\) test norm. The inner limit exists and equals
\(\Rcal_h(I_hg_0)\). Refining that interpolant gives
\(\Ecal(g_0)\) by \eqref{eq:uniform}. At fixed \(\varepsilon\),
\eqref{eq:res} gives zero as \(h\downarrow0\). Thus the two
iterated assertions follow as well as both simultaneous regimes.
\end{proof}

\section{A comparable-scale quadrature resonance}\label{sec:alias}

The rational tilt and its fixed origin permit exact cancellation
of the leading interpolated cosine, not just a small-wavelength
estimate. Write
\(\dd u\,L(\mathbb Z^4)\subset q\mathbb Z\) for a fixed positive
lattice spacing \(q\). On the common tilted family, take
\begin{equation}
                        \varepsilon=\frac{q h}{2\pi}.
 \label{eq:critical}
\end{equation}
Then \(u/\varepsilon\) is an integer multiple of \(2\pi\)
at every mesh vertex. In particular the initial phase is fixed
at the mesh origin, rather than chosen independently from cell
to cell. By the tilt choice, endpoints of every template edge
have distinct \(u\); the same is true of all pairs of vertices
on a triangle.

Use the constant-envelope family with either
\(w_\varepsilon=\kappa_0\varepsilon\cos(u/\varepsilon)\) or
\(w_\varepsilon=\kappa_0\varepsilon\sin(u/\varepsilon)\), and
the initial data of \eqref{eq:vac} at zero. Both limits have
\(a_0=\cos(\kappa_0 u/\sqrt2)\). In fact
\(a_\varepsilon-a_0=O(\varepsilon^2)\) in \(C^0\).
The constant-oscillator formula has the kernel and forcing
\[
    Q(u-s)=\frac{\sin(\kappa_0(u-s)/\sqrt2)}{\kappa_0/\sqrt2},
    \qquad
    \pm(\kappa_0^2/2)\cos(2s/\varepsilon)a_\varepsilon(s).
\]
The first integration by parts gives zero endpoints since
\(\sin(2s/\varepsilon)\) vanishes at \(s=0\) and \(Q(0)=0\)
at the other endpoint; a second gives the asserted
\(O(\varepsilon^2)\) using bounded \(C^2\) norms of
\(a_\varepsilon\). This is valid on either side of \(u=0\).
The leading cosine perturbation is therefore
\[
     2\kappa_0\varepsilon a_0(u)^2\cos(u/\varepsilon)S_1,
     \qquad S_1=\dd x^2-\dd y^2 ,
\]
with uniform \(O(\varepsilon^2)\) error in metric values.

Freeze the smooth coefficient \(a_0^2\) per simplex. Each edge
sees a nonzero integer phase increment \(2\pi m\), so both
\(\int_0^1\cos(2\pi m t)\dd t\) and
\(\int_0^1t\cos(2\pi m t)\dd t\) vanish. On a reference triangle
the constant complex moment is
\[
 \int_0^1\int_0^{1-s}e^{i(ps+q' t)}\,\dd t\,\dd s
    =\frac{1}{i q'}\left[
       e^{i q'}\int_0^1 e^{i(p-q')s}\,\dd s
          -\int_0^1e^{ips}\,\dd s\right]=0,
\]
where \(p,q',p-q'\) are nonzero integer multiples of \(2\pi\).
A constant tangential tensor contraction just multiplies this
integral by a constant. Therefore the frozen cosine perturbation
has zero for \emph{all} interpolation moments, not only the
first moments of a particular component. The freezing error
is \(O(h\varepsilon)\), and \eqref{eq:interp} gives
\[
 \|I_h(g_\varepsilon^{\cos}-g_0)\|_{\infty,h}=O(h^2),
 \qquad
 \|D_h I_h(g_\varepsilon^{\cos}-g_0)\|_{\infty,h}=O(h).
\]
Inserting in \eqref{eq:uniform} about \(g_0\) proves
\(\Rcal_h(I_hg_\varepsilon^{\cos})=\Ecal(g_0)+O(h^2)\)
in \(H^{-2}\), even though \(\varepsilon/h\) is a fixed positive
constant. Its centered broken derivative tends strongly to zero.

For the sine on the very same mesh and ratio, a repeating
near-\(x\) edge can be kept with \(S_1(v_e,v_e)\ne0\) for
its auxiliary unit tangent \(v_e\). A nonzero phase increment
integer \(m_e\) gives zero mean but nonzero first moment
\(\int_0^1 t\sin(2\pi m_e t)\dd t=-1/(2\pi m_e)\).
On an affine edge restriction the centered first moment equals
one twelfth of its slope in the parameter \(t\). Its directional
gradient in the coordinates is therefore bounded away from
zero after the leading \(\varepsilon\) amplitude is divided
by the physical edge length comparable to \(h\); the
\(O(h^2)\) metric errors contribute only \(O(h)\).
A positive periodic fraction of cells contains this edge
direction. Since their affine derivative tensor is constant,
the \(L^2\) norm of the sine's centered broken gradient has
positive liminf on any meshed compact interior box of
positive volume. This last statement is about retained
derivative signal, not a claim that the sine Einstein residual
has a different limit.

\section{A circular-polarization collision}\label{sec:helicity}

Let \(S_1=\operatorname{diag}(1,-1)\) and
\(S_2=\left(\begin{smallmatrix}0&1\\1&0\end{smallmatrix}\right)\)
on the marked transverse screen, and \(\theta=u/\varepsilon\).
Define two transverse Rosen matrices by
\begin{align}
 \gamma_{\varepsilon,\pm}
   &=a_\varepsilon^2 C_{\varepsilon,\pm},&
 C_{\varepsilon,\pm}
   &=\exp\!\left(
       2\kappa_0\varepsilon(\cos\theta\,S_1\pm\sin\theta\,S_2)
              \right),\notag\\
 a_\varepsilon
   &=\cos\!\left(
       \frac{\sinh(2\kappa_0\varepsilon)}{2\varepsilon}u
          \right).
 \label{eq:circular}
\end{align}
The Lorentzian metric has the same \(-2\,\dd u\,\dd v\) block.
For \(P_\pm=\cos\theta\,S_1\pm\sin\theta\,S_2\),
\(P_\pm^2=I\), \(P_\pm\dot P_\pm=-\dot P_\pm P_\pm\), and
\(\dot P_\pm^2=I\), where dots are \(\theta\)-derivatives.
Writing \(C_{\varepsilon,\pm}=c_\varepsilon I+s_\varepsilon P_\pm\)
with hyperbolic cosine and sine, one obtains
\[
       (C_{\varepsilon,\pm}^{-1}C_{\varepsilon,\pm}')^2
               =\frac{s_\varepsilon^2}{\varepsilon^2}I .
\]
For \(\det C=1\) the Rosen Ricci formula is
\(\Ric_{uu}=-2a''/a-\tr[(C^{-1}C')^2]/4\), directly from the
transverse Christoffel contraction. This proves that both waves
in \eqref{eq:circular} are exactly vacuum. They have the same
limit \(\bar a=\cos(\kappa_0 u)\) and
\(\bar G_{uu}=2\kappa_0^2\). Also
\(a_\varepsilon-\bar a=O(\varepsilon^2)\) in \(C^1\), and the
bounds \eqref{eq:jets} and the admissible ball still apply
for small \(\kappa_0\).

In the common normalized screen the continuum derivative
has coefficients
\[
       2\kappa_0(-\sin\theta,\ \pm\cos\theta)
       \qquad\text{in the basis }(S_1,S_2).
\]
At positive \(u\)-frequency this gives, up to the same
positive weight, the Hermitian outer product of
\((i,\pm1)\). The two matrices are orthogonal rank-one
helicity polarizations with equal trace and opposite
circular off-diagonal part. The negative frequency has the
conjugate matrix. This is a comparison in the same marked
orientation, not the assertion that a bare coordinate
component without its frame is invariant. Vacuum examples already show why
polarization matters even at fixed scalar trace~\cite{VacuumPolarization}.

In the resolved regime the broken metric gradient of each
interpolant differs in \(L^2\) from its signed continuum gradient
by \(O(h/\varepsilon)\), so fixed bounded quadratic microlocal
tests retain these distinct screen tensors. In the underresolved
regime both gradients converge strongly to the smooth limiting
gradient by \eqref{eq:interp} and \eqref{eq:jets}, even though
\eqref{eq:under} retains the scalar dust.

\begin{theorem}[Critical alias collision]\label{thm:circular}
Along the particular critical diagonal \eqref{eq:critical},
the bounded quadratic broken-gradient tests of the two centered
interpolated metrics \eqref{eq:circular} have vanishing difference.
This includes order-zero microlocal tests in a fixed marked
screen. The corresponding signed Einstein functionals satisfy
\[
 \|\Rcal_h(I_hg_{\varepsilon,+})
          -\Rcal_h(I_hg_{\varepsilon,-})\|_{H^{-2}}=O(h).
\]
This equality of leading outputs does not require assigning
a limiting Einstein density to either critical wave separately.
\end{theorem}

\begin{proof}
The common cosine quadrature is erased by exactly the
same moment calculation as in Section~\ref{sec:alias}. Let
\[
       I_hg_{\varepsilon,\pm}=\bar g+\delta_h\pm d_h,
 \qquad \bar g=-2\,\dd u\,\dd v+\bar a^2(\dd x^2+\dd y^2).
\]
Here \(\delta_h\) includes \(I_h\bar g-\bar g\). The common
cosine error after freezing is \(O(h\varepsilon)\), and the
even terms \(a_\varepsilon-\bar a\) and
\(c_\varepsilon-1\) cost \(O(\varepsilon^2)\). The odd
sine perturbation has size \(O(\varepsilon)\). Thus the
affine estimate gives
\begin{align*}
 \|\delta_h\|_{\infty,h}&=O(h^2),&
 \|D_h\delta_h\|_{\infty,h}&=O(h),&
 D_h^2\delta_h&=-D^2\bar g=O(1),\\
 \|d_h\|_{\infty,h}&=O(h),&
 \|D_h d_h\|_{\infty,h}&=O(1),&
 D_h^2 d_h&=0 .
\end{align*}
The two centered gradients are therefore negatives of each
other modulo \(O(h)\) in \(L^2\). If the centering subtracts
their corresponding weak subsequential limits instead, the
limits themselves are opposite, and the same statement
persists. Pairing a bounded \(L^2\) quadratic operator with
these two gradients, Cauchy--Schwarz shows that the difference
is \(O(h)\).

For the signed Einstein terms, expand around \(\bar g\)
through quadratic order using the same branch and polynomial
jet structure as in \eqref{eq:uniform}. The linear contribution
is \(O(h)\) for the difference, by the transferred smooth
Hessian \eqref{eq:hess} evaluated on \(2d_h\). The pure
quadratic squares of \(d_h\) cancel between \(+\) and \(-\).
Mixed quadratic terms pairing \(d_h\) with \(\delta_h\)
are \(O(h)\) in the bulk, \(O(h^2)\) per face area and
\(O(h^3)\) per hinge area. In particular the bulk
\(d_h D_h^2\delta_h\) term is still \(O(h)\), although that
second jet is only bounded. The leading cubic remainders have
size \(O(h)\) in the bulk (each term with two first jets
contains an \(O(h)\) metric factor), \(O(h^2)\) per face,
and \(O(h^3)\) per hinge. This is a Taylor estimate of the
smooth coefficient functions, not a claim that the first
jet error itself is small at critical scale. The same trace
sums \(h^{-1},h^{-2}\) give the claimed \(H^{-2}\) bound.
\end{proof}

\section{What the noncommutation establishes}

For the exact waves the smooth vacuum Einstein density is zero
at every wavelength. The coarse effective stress is instead the
Einstein density of the uniform limiting metric. The two
limits are realized by the \emph{same} signed Lorentzian
Regge interpolation on a uniform non-null branch. The contrast
does not require singular causal angles or a different
triangulation tuned separately to every smooth metric.
The critical examples also show why scalar effective stress and
a nominal finite resolution do not by themselves determine
which radiative polarization data an interpolant retains:
the actual affine moment response can annihilate one quadrature.

The statements concern this degree-one metric interpolation and
the complete weak variational Einstein tensor. They do not assert
discrete solution existence or convergence, convergence under
edge-length-only sampling, behavior near null hinges, a limiting
quantum amplitude, or a universal law at every finite scale ratio
\(h/\varepsilon\). The alias collision concerns quadratic
readouts of the broken gradients and the weak Einstein residual;
the raw finite-mesh coefficients need not be identical.

\paragraph{Working provenance.} This manuscript was prepared with AI assistance and internally cross-checked. No external peer review is asserted.

\begin{thebibliography}{9}
\bibitem{SignedRegge}
\emph{A signed Lorentzian Regge approximation theorem for the
Einstein tensor}, paper 28, AI-Automated Quantum Gravity Research,
\href{https://github.com/g3173028019-beep/ai-automated-quantum-gravity-research/tree/main/papers/28-lorentzian-regge-einstein-tensor}
{source and manuscript}.

\bibitem{Burnett}
G.~A.~Burnett, \emph{The high-frequency limit in general relativity},
J.~Math.~Phys.~\textbf{30} (1989), 90--96,
\href{https://doi.org/10.1063/1.528594}{doi:10.1063/1.528594}.

\bibitem{GreenWald}
S.~R.~Green and R.~M.~Wald,
\emph{Examples of backreaction of small-scale inhomogeneities
in cosmology}, Phys.~Rev.~D \textbf{87} (2013), 124037,
\href{https://arxiv.org/abs/1304.2318}{arXiv:1304.2318}.

\bibitem{HuneauLuk}
C.~Huneau and J.~Luk,
\emph{Burnett's conjecture in generalized wave coordinates},
\href{https://arxiv.org/abs/2403.03470}{arXiv:2403.03470}.

\bibitem{Hayward}
G.~Hayward, \emph{Gravitational action for spacetimes with
nonsmooth boundaries}, Phys.~Rev.~D \textbf{47} (1993), 3275--3280,
\href{https://doi.org/10.1103/PhysRevD.47.3275}
{doi:10.1103/PhysRevD.47.3275}.

\bibitem{Sorkin}
R.~D.~Sorkin, \emph{Lorentzian angles and trigonometry including
lightlike vectors}, \href{https://arxiv.org/abs/1908.10022}
{arXiv:1908.10022}.

\bibitem{Asante}
S.~K.~Asante, B.~Dittrich and J.~Padua-Arg\"uelles,
\emph{Complex actions and causality violations: Applications to
Lorentzian quantum cosmology},
\href{https://arxiv.org/abs/2112.15387}{arXiv:2112.15387}.

\bibitem{GawlikNeunteufel}
E.~S.~Gawlik and M.~Neunteufel,
\emph{Finite element approximation of the Einstein tensor},
\href{https://arxiv.org/abs/2310.18802}{arXiv:2310.18802}.

\bibitem{VacuumPolarization}
\emph{Physical polarization in the high-frequency limit of vacuum
gravity}, paper 29B, AI-Automated Quantum Gravity Research,
\href{https://github.com/g3173028019-beep/ai-automated-quantum-gravity-research/tree/main/papers/29b-high-frequency-graviton-polarization}
{source and manuscript}.
\end{thebibliography}

\end{document}
